% 乔治·伽莫夫（综述）
% license CCBYSA3
% type Wiki

本文根据 CC-BY-SA 协议转载翻译自维基百科\href{https://en.wikipedia.org/wiki/George_Gamow}{相关文章}

乔治·伽莫夫（George Gamow，有时写作 Gammoff；原名 Гео́ргий Анто́нович Га́мов，Georgiy Antonovich Gamov，1904 年 3 月 4 日－1968 年 8 月 19 日）是一位苏联和美国的博学家、理论物理学家与宇宙学家。他是乔治·勒梅特大爆炸理论的早期倡导者与发展者之一。伽莫夫首先通过量子隧穿理论解释了α衰变的机理，提出了液滴模型——第一个原子核的数学模型，并研究了放射性衰变、恒星形成、恒星核合成、以及大爆炸核合成（他统称为“核宇宙成因论”，nucleocosmogenesis），并预言了宇宙微波背景辐射的存在以及分子遗传学的基础。伽莫夫是量子隧穿现象发展与理解的关键人物之一。

在中晚期生涯中，伽莫夫将大量精力投入教学，并撰写了多部科普著作，包括《一二三……无穷大》以及“汤普金斯先生”系列丛书（1939–1967）。他的一些著作在最初出版半个多世纪后仍在印行。为纪念他，美国科罗拉多大学博尔德分校设立了乔治·伽莫夫纪念讲座。
\subsection{早年生平与职业生涯}
伽莫夫出生于俄国帝国敖德萨（今乌克兰敖德萨）。他的父亲在中学教授俄语与文学，母亲则在女子学校教授地理与历史。除了俄语外，伽莫夫还从母亲那里学会了一些法语，并由家庭教师学习了德语。在大学期间，他又学习了英语，并最终能流利使用。伽莫夫早期的大部分论文以德语或俄语撰写，后来则以英语发表技术论文和科普著作。

他先后就读于敖德萨物理与数学学院（1922–1923），以及列宁格勒大学（1923–1929）。在列宁格勒期间，伽莫夫师从亚历山大·弗里德曼，直到弗里德曼于1925年早逝，被迫更换导师。在大学时期，伽莫夫结识了三位理论物理专业的同学——列夫·朗道、德米特里·伊万年科和马特维·布龙施泰因。四人组成了一个自称为“三剑客”的小组，定期聚会讨论并分析当时发表的开创性量子力学论文。后来，伽莫夫也用同样的称呼来形容他与拉尔夫·阿尔费和罗伯特·赫尔曼共同组成的研究小组。
